% ---
% Reemplaza líneas 2577–2629 de Docs/Technical_Report.tex.
% Decisiones que lo respaldan: D-001, D-012, H-001, H-002, H-010.
% ---

\subsection{Módulo 1: Lectura de datos}
\label{sec:modulo1}

El primer módulo gestiona la adquisición y unificación de
los datos del dataset CBIS-DDSM, disponible públicamente
en The Cancer Imaging Archive (TCIA). El dataset comprende 1,566
pacientes con imágenes mamográficas digitalizadas en
formato DICOM y cuatro archivos CSV con metadatos clínicos
y anotaciones patológicas validadas.

La lectura de los archivos DICOM se realiza mediante la
librería \texttt{pydicom}, extrayendo el arreglo de píxeles
en formato de 16 bits y los metadatos clínicos relevantes:
\texttt{PatientID}, \texttt{Modality},
\texttt{Rows}, \texttt{Columns}, \texttt{BitsStored} y
\texttt{PhotometricInterpretation}. El análisis exploratorio
confirmó que la totalidad de las imágenes del dataset
utilizan la convención fotométrica \texttt{MONOCHROME2},
donde los valores altos corresponden al blanco (tejido
denso) y los valores bajos al negro (fondo), lo que
elimina la necesidad de inversión de píxeles en el
preprocesamiento.

Ninguna de las 3\,103 mamografías incluye las etiquetas \texttt{PixelSpacing},
\texttt{ImagerPixelSpacing} ni \texttt{Manufacturer}. Sin un espaciado físico, las características de
forma se expresan en píxeles, y no es posible distinguir el digitalizador de origen.

Los cuatro archivos CSV del dataset
(\texttt{mass\_case\_description\_train\_set},
\texttt{mass\_case\_description\_test\_set},
\texttt{calc\_case\_description\_train\_set},
\texttt{calc\_case\_description\_test\_set})
se unifican en un DataFrame mediante
\texttt{pandas}, agregando las columnas
\texttt{finding\_type} (masa / calcificación) y
\texttt{split} (train / test). La etiqueta objetivo
\texttt{pathology} se normaliza a una variable binaria:
\texttt{BENIGN\_WITHOUT\_CALLBACK} se trata equivalente
a \texttt{BENIGN}, resultando en la codificación
0\,=\,benigno y 1\,=\,maligno.

El enlace entre los registros CSV y los archivos DICOM
en disco se resuelve mediante un índice construido
recursivamente sobre la estructura de carpetas del
NBIA Data Retriever, identificando para cada caso el
archivo de mamografía completa, la máscara ROI binaria
y la imagen recortada. La clasificación entre máscara
y recorte se realiza leyendo el contenido de píxeles:
las máscaras presentan exactamente dos valores únicos
(0 y 255) mientras que las imágenes recortadas presentan
cientos de valores distintos.

La descarga inicial resultó incompleta, lo que dejaba la cobertura en el 97.6\,\%.
Las 84 series faltantes se recuperaron mediante la API REST de TCIA, solicitando
cada serie por su identificador, después de que el NBIA Data Retriever rechazara
un manifiesto reconstruido. Con ello la cobertura alcanza los 3\,568 registros.
Estos corresponden a 3\,103 mamografías distintas, porque varias lesiones
comparten la misma imagen.

La salida de este módulo es un DataFrame unificado con
3,568 registros, con rutas resueltas a los archivos DICOM
correspondientes y etiquetas binarias validadas, que sirve
como base para todos los módulos posteriores.
